

\begin{figure}[htbp]
\centering
\begin{tikzpicture}[
    scale=0.85,
    >=Latex,
    line cap=round,
    line join=round,
    every node/.style={font=\small},
    lattice/.style={circle, fill=black, inner sep=1.15pt},
    feasiblept/.style={circle, fill=black, inner sep=1.55pt},
    lpregion/.style={draw=blue!65!black, fill=blue!8, line width=0.0pt},
    cutregion/.style={draw=red!70!black, fill=red!12, line width=0.0pt},
    constraint/.style={blue!70!black, line width=1.0pt},
    cut/.style={red!75!black, dashed, line width=1.1pt},
    solLP/.style={circle, fill=blue!80!black, inner sep=2pt},
    solIP/.style={circle, fill=green!55!black, inner sep=2pt}
]

% --------------------------------------------------
% Parameters
% --------------------------------------------------
\def\eps{0.20}

% Constraint 1: x_2 <= a_1 x_1 + b_1
\pgfmathsetmacro{\aone}{(-1 - 2*\eps)/4}
\pgfmathsetmacro{\bone}{3 + \eps}

% Constraint 2: x_2 <= a_2 x_1 + b_2,
% written as x_2 <= a_2 (x_1 - (5 + eps))
\pgfmathsetmacro{\atwo}{(-3 + 2*\eps)/2}
\pgfmathsetmacro{\btwo}{-\atwo*(5+\eps)}

% Intersection of the two LP constraints = fractional LP optimum
\pgfmathsetmacro{\xLP}{(\btwo-\bone)/(\aone-\atwo)}
\pgfmathsetmacro{\yLP}{\aone*\xLP+\bone}

% Values used for plotting boundaries
\pgfmathsetmacro{\yA}{\bone}
\pgfmathsetmacro{\yB}{\aone*5+\bone}
\pgfmathsetmacro{\yC}{\atwo*5+\btwo}
\pgfmathsetmacro{\xD}{-\btwo/\atwo}

% A Gomory-like illustrative cut, chosen to remove x^LP
% while preserving the integer points below it.
\pgfmathsetmacro{\acut}{-1}
\pgfmathsetmacro{\bcut}{5}

% Intersections of cut with the two LP edges
\pgfmathsetmacro{\xCutOne}{(\bcut-\bone)/(\aone-\acut)}
\pgfmathsetmacro{\yCutOne}{\acut*\xCutOne+\bcut}

\pgfmathsetmacro{\xCutTwo}{(\btwo-\bcut)/(\acut-\atwo)}
\pgfmathsetmacro{\yCutTwo}{\acut*\xCutTwo+\bcut}

% --------------------------------------------------
% Axes
% --------------------------------------------------
\draw[->, thick] (-0.25,0) -- (5.75,0) node[below right] {$x_1$};
\draw[->, thick] (0,-0.25) -- (0,4.35) node[above left] {$x_2$};

% --------------------------------------------------
% Lattice points
% --------------------------------------------------
\foreach \x in {0,...,5}{
  \foreach \y in {0,...,4}{
    \node[lattice, fill=black!35] at (\x,\y) {};
  }
}

% --------------------------------------------------
% Original LP relaxation
% --------------------------------------------------
\filldraw[lpregion]
  (0,0) --
  (\xD,0) --
  %(5,\yC) --
  (\xLP,\yLP) --
  (0,\yA) --
  cycle;


% --------------------------------------------------
% Region removed by the cut
% --------------------------------------------------
\filldraw[cutregion]
  (\xCutOne,\yCutOne) --
  (\xLP,\yLP) --
  (\xD,0) --
  (5,0) --
  cycle;


% Red hatching inside removed part
\begin{scope}
  \clip
    (\xCutOne,\yCutOne) --
    (\xLP,\yLP) --
    (\xD,0) --
    (5,0) --
    cycle;

  \foreach \s in {-2,-1.75,...,6}{
    \draw[red!65, line width=0.65pt]
      (\s,0) -- ++(3.0,3.0);
  }
\end{scope}
% --------------------------------------------------
% Constraint boundaries
% --------------------------------------------------
\draw[constraint] (0,\yA) -- (\xLP,\yLP);
\draw[constraint] (\xLP,\yLP) -- (5,\yC);

% Optional extension of LP boundary lines
\draw[constraint, opacity=0.6]
  (-0.15,{\aone*(-0.15)+\bone})
  --
  (5.45,{\aone*5.45+\bone});

% Stop the second boundary at x_2=0 to avoid drawing below the axis
\draw[constraint, opacity=0.6]
  (2.65,{\atwo*2.65+\btwo})
  --
  ({-\btwo/\atwo},0);

% --------------------------------------------------
% Cutting plane
% --------------------------------------------------
\draw[cut]
  (1.0,4.0)
  --
  (5,0);

\node[red!75!black, above right] at (1.55,{\acut*1.55+\bcut}) {cut};

% --------------------------------------------------
% Feasible integer points: those satisfying both LP constraints and the cut
% --------------------------------------------------
\foreach \x in {0,...,5}{
  \foreach \y in {0,...,4}{
    \pgfmathtruncatemacro{\ok}{
      ifthenelse(
        \y <= \aone*\x+\bone + 0.001 &&
        \y <= \atwo*\x+\btwo + 0.001 &&
        \y <= \acut*\x+\bcut + 0.001,
        1,0)
    }
    \ifnum\ok=1
      \node[feasiblept] at (\x,\y) {};
    \fi
  }
}

% --------------------------------------------------
% Solutions
% --------------------------------------------------
\node[solLP] at (\xLP,\yLP) {};
\node[blue!80!black, above right] at (\xLP,\yLP) {$x^{\mathrm{LP}}$};

\node[solIP] at (3,2) {};
\node[green!45!black, below left] at (3,2) {$x^{\mathrm{IP}}$};

% --------------------------------------------------
% Small label for feasible region
% --------------------------------------------------
%\node[blue!65!black] at (1.25,1.35) {$P_{\mathrm{LP}}$};

\end{tikzpicture}
\caption{Illustration of a cutting plane. The blue shaded region represents the LP relaxation, whose optimum $x^{\mathrm{LP}}$ is attained at a fractional vertex. The dashed red inequality cuts off this fractional solution while preserving the integer-feasible points, including $x^{\mathrm{IP}}$.}
\label{fig:cutting_plane}
\end{figure}